\section{Quota-Constrained Multi-Rate Scheduling}
\label{sec:method}

The scheduler separates \emph{what} is sent (tiers), \emph{when} (per-tier intervals chosen under the quota), \emph{how large} (an optional compressor) and \emph{how private} (the DP mechanism and its accountant). Fig.~\ref{fig:arch} shows the architecture. The server-side controller observes the received slices, the sites observe only the broadcast, and the link and the privacy mechanism sit between them.

\begin{figure*}[t]
\centering
\begin{tikzpicture}[x=1mm,y=1mm]
\node[hdr] at (0,8) {Server-side controller (\tierfed{})};
\node[hdr] at (62,8) {Link and policy};
\node[hdr] at (124,8) {Site $k$ of $N$};
% server column
\node[srvbox, text width=42mm] (sched) at (0,0)   {\textbf{Schedule.} Tiers due $\mathcal A_r$ by (\ref{eq:due})};
\node[srvbox, text width=42mm] (solve) at (0,-12) {\textbf{Exact solver} (\ref{eq:ip}). $\mathbf K$ s.t.\ $\sum_g w_g/K_g\le\rho$};
\node[srvbox, text width=42mm] (disp)  at (0,-24) {\textbf{Dispersion proxy} $\Gh_g$ (\ref{eq:gammahat}) from the received slices};
\node[srvbox, text width=42mm] (agg)   at (0,-36) {\textbf{Aggregation} (\ref{eq:agg}). Under DP: clip, sum, add noise, fixed denominator (\ref{eq:dp})};
% link column
\node[lnkbox, text width=32mm] (op)   at (62,-8)  {\textbf{Operator.} Sets $\rho$; under DP also $\varepsilon$, $\delta$, $q$};
\node[lnkbox, text width=32mm] (acc)  at (62,-22) {\textbf{Accountant} (\ref{eq:rdp}). Calibrates $z$ to the planned releases};
\node[lnkbox, text width=32mm] (link) at (62,-36) {\textbf{Uplink} $(\beta,\ell,P)$. Compressor with error feedback};
% site column
\node[stebox, text width=42mm] (bcast) at (124,0)   {Receive the active tiers; restore the private head $\theta_{f,k}$};
\node[stebox, text width=42mm] (train) at (124,-12) {One local epoch (Adam), $s_k=\lceil n_k/B\rceil$ steps};
\node[stebox, text width=42mm] (ret)   at (124,-24) {Return the active tiers as an update w.r.t.\ the broadcast};
\node[stebox, text width=42mm] (eval)  at (124,-36) {Select on the site's validation split, report on its test split};
\begin{scope}[on background layer]
\node[grp, fit=(sched)(solve)(disp)(agg)] {};
\node[grp, fit=(op)(acc)(link)] {};
\node[grp, fit=(bcast)(train)(ret)(eval)] {};
\end{scope}
% data path down the round (solid)
\draw[arr] (sched.east) -- node[lab, above] {broadcast $\{\theta_g\}_{g\in\mathcal A_r}$} (bcast.west);
\draw[arr] (bcast.south) -- (train.north);
\draw[arr] (train.south) -- (ret.north);
\draw[arr] (ret.west) -- ++(-9,0) node[lab, above, pos=0.5] {$b_k$ bits} |- (link.east);
\draw[arr] (link.west) -- node[lab, above] {$\sum_k$} ([yshift=-2.5mm]agg.east);
% policy inputs (solid)
\draw[arr] (op.west) -- ++(-11,0) node[lab, above, pos=0.5] {$\rho$} |- (solve.east);
\draw[arr] (op.south) -- node[lab, right] {$\varepsilon,\delta,q$} (acc.north);
\draw[arr] (acc.west) -- ++(-7,0) node[lab, above, pos=0.5] {$z$} |- ([yshift=2.5mm]agg.east);
% measurements fed back (dashed)
\draw[meas] (agg.north) -- node[lab, right] {received slices} (disp.south);
\draw[meas] (disp.north) -- node[lab, right] {$\Gh_g$} (solve.south);
\draw[arr]  (solve.north) -- node[lab, right] {$\mathbf K$} (sched.south);
\draw[meas] (train.east) -- ++(4,0) |- (eval.east);
\end{tikzpicture}
\caption{Scheduler and environment. Solid arrows carry data and policy inputs; dashed arrows are measurements fed back. Under DP the dashed path from the received slices to the solver is cut and the schedule is planned before training (Section~\ref{sec:method-privacy}).}
\label{fig:arch}
\end{figure*}

\subsection{Dispersion Measurement}
\tierfed{} estimates the heterogeneity of each tier from the dispersion of the slices it receives. For tier $g$ with tensors $j\in g$,
\begin{equation}
\Gh_g^{(r)}=\frac{\sum_{j\in g}\frac1{|\mathcal S_r|}\sum_{k\in\mathcal S_r}\big\|\theta_{k,j}^{(r)}-\bar\theta_j^{(r)}\big\|_2^2}{\sum_{j\in g}\big\|\bar\theta_j^{(r)}\big\|_2^2+\xi_0},
\label{eq:gammahat}
\end{equation}
a scale-free ratio ($\xi_0=10^{-12}$) that makes tiers of different magnitude comparable. It exists only for tiers received in round $r$; otherwise the last value is carried forward. The received slices result from $m_g=(r-t_g)\,s$ local steps since the tier's last synchronisation $t_g$, and the standard drift lemma of local SGD~\cite{stich2019,khaled2020} bounds the expected numerator by a term of order $\eta^2(m_g\sigma_g^2+m_g^2\Gamma_g)$. $\Gh_g$ can therefore grow quadratically with the interval that produced it, so the controller partly measures its own decisions. The debiased proxy $\tilde\Gamma_g=\Gh_g/(r-t_g)^2$ removes the quadratic growth of the heterogeneity term; it is released as an option but not evaluated here.

\subsection{The Programme and Its Closed Form}
\emph{Bound.} The analysis concerns the shared backbone with the private heads frozen; the personalisation gap this introduces is measured empirically. Assume block-separable $L$-smoothness (the gradient of tier $g$ is $L$-Lipschitz in $\theta_g$ and does not depend on the other tiers; this ignores cross-tier coupling and is the main idealisation), unbiased stochastic gradients with per-tier variance $\sigma_g^2$, bounded per-tier heterogeneity $\frac1N\sum_k\|\nabla_gF_k(\theta)-\nabla_gF(\theta)\|^2\le\Gamma_g$, $F\ge F^\star$~\cite{stich2019,wang2021coop}, and the same number $s$ of local steps at every site. With unequal $s_k$, plain averaging converges towards a reweighted objective~\cite{wang2020fednova}; the experiments keep $s_k=\lceil n_k/B\rceil$ and do not normalise.
\begin{proposition}[Multi-rate local-SGD bound]
\label{prop:bound}
With local SGD of step size $\eta$ and tier $g$ averaged every $K_g$ rounds, there are absolute constants $c_0,c_1,c_2$ such that for $\eta\le c_0/(LK_{\max}s)$,
\begin{align}
\frac1R\sum_{r=1}^R\mathbb E\big\|\nabla F(\bar\theta_r)\big\|^2
&\le \frac{c_1(F(\bar\theta_1)-F^\star)}{\eta sR}+c_1\eta L\sum_g\frac{\sigma_g^2}{N}\notag\\
&\quad+c_2\,\eta^2L^2s^2\sum_g K_g^2\big(\Gamma_g+\sigma_g^2\big).
\label{eq:bound}
\end{align}
\end{proposition}
The proof applies the perturbed-iterate argument of~\cite{wang2021coop} blockwise (supplement, Section~S2). In the bound, a long interval is cheap only for a tier with small $\Gamma_g+\sigma_g^2$.

\emph{Programme.} Given a quota $\rho$ the controller solves
\begin{equation}
\mathbf K^\star=\arg\min_{K_g\in\{1,\dots,K_{\max}\}}\sum_{g\in\mathcal G_{\mathrm{sh}}}K_g^2\,\Gh_g\ \ \text{s.t.}\ \sum_{g\in\mathcal G_{\mathrm{sh}}}\frac{w_g}{K_g}\le\rho .
\label{eq:ip}
\end{equation}
It keeps the $K_g^2\Gamma_g$ term of (\ref{eq:bound}). The $K_g^2\sigma_g^2$ term is not separately estimable from parameter dispersion, which mixes it with the heterogeneity term, and the common factor $s^2$ does not change the minimiser. With the head private, $\mathcal G_{\mathrm{sh}}=\{\mathrm s,\mathrm m\}$ and $K_{\max}=8$, so the programme has 64 candidates and is solved by enumeration, with ties broken towards the denser schedule.

\begin{proposition}[Closed form and the value of differentiation]
\label{prop:cuberoot}
Let $\Gamma_g>0$, $w_g>0$, $\sum_gw_g=1$, and relax $K_g$ to the positive reals. Then (\ref{eq:ip}) has the unique minimiser
\begin{equation}
K_g^\star=\Big(\frac{w_g}{\Gamma_g}\Big)^{1/3}\frac{S}{\rho},\qquad S=\sum_hw_h^{2/3}\Gamma_h^{1/3},
\label{eq:cuberoot}
\end{equation}
with value $J^\star=S^3/\rho^2$. The uniform interval $K=1/\rho$ costs $J_{\mathrm u}=\sum_g\Gamma_g/\rho^2$, and
\begin{equation}
\Psi=1-\frac{J^\star}{J_{\mathrm u}}=1-\frac{S^3}{\sum_g\Gamma_g}\in[0,1),
\label{eq:psi}
\end{equation}
with $\Psi=0$ if and only if $w_g/\Gamma_g$ is equal for all tiers. When $1/\rho$ is an integer, no integer schedule lowers the objective of (\ref{eq:ip}) below that of the uniform period by more than the fraction $\Psi$.
\end{proposition}
The objective is strictly convex and increasing in every $K_g$ and the feasible set is convex, so the quota binds and the KKT point is the minimiser; $S^3\le(\sum_gw_g)^2\sum_g\Gamma_g$ is H\"older's inequality with exponents $3/2$ and $3$; the integer statement holds because the relaxation bounds the integer optimum from below and $K=1/\rho$ is integer-feasible (supplement, Section~S2). Two consequences shape the evaluation. The interval ratio $K_a/K_b=(w_a\Gamma_b/w_b\Gamma_a)^{1/3}$ compresses differences: a factor of eight in payload-to-dispersion ratio is needed for a factor of two in interval, so tiers with similar ratios collapse onto one interval after rounding. And $\Psi$ is scale-free in $\Gamma$, so the common bias of $\Gh_g$ cancels when it is measured at the end of warm-up, where every tier has been synchronised every round; it then predicts, before any sparse round, how much a per-tier schedule can gain over a uniform one.

\subsection{Schedule Rule, Phases and the Quota}
Rounds are indexed $r=1,\dots,R$. The first $T_w$ rounds synchronise every shared tier, and thereafter tier $g$ is due when
\begin{equation}
\mathcal A_r=\begin{cases}\mathcal G_{\mathrm{sh}}, & r\le T_w,\\ \{g:\ (r-T_w-\phi_g)\bmod K_g=0\}, & r>T_w,\end{cases}
\label{eq:due}
\end{equation}
with phase $\phi_g=0$ unless staggered. The fixed periods, the DP planner and the accountant all use this one rule, and the consistency check of Section~\ref{sec:results-check} recomputes every release from it. Adaptive \tierfed{} sends tier $g$ when $r-t_g\ge K_g$, which coincides with (\ref{eq:due}) whenever $\mathbf K$ is constant after warm-up. $T_w=5$ and $K_{\max}=8$ are fixed across every dataset, architecture, loss and privacy setting. Fig.~\ref{fig:round} shows one server round.

\emph{Phases and the peak.} A mean quota does not bound the payload $P_r$ of a single round: aligned tiers coincide and send the full shared payload. Two tiers with phases $\phi_a,\phi_b$ are never due in the same round if and only if $\phi_a-\phi_b\not\equiv0\pmod{\gcd(K_a,K_b)}$ (Chinese remainder theorem), so with $\gcd\ge2$ the offset $\lfloor\gcd/2\rfloor$ separates them and bounds the post-warm-up peak by $\max_gw_g$ at an unchanged mean; with co-prime intervals no phase helps. The staggered variant adds $\max_{r>T_w}P_r\le\pi$ to (\ref{eq:ip}) and checks it over one hyper-period $\mathrm{lcm}(\mathbf K)$. A full-rate warm-up still sends the whole payload, so the warm-up rounds can be sent at eight bits.

\emph{The quota.} Solving (\ref{eq:ip}) with $\rho$ bounds $\Pi_\infty$ of (\ref{eq:payload}), not the realised $\Pi_R$: a warm-up of $T_w$ rounds at payload fraction $c_w$ adds $T_w(c_w-\Pi_\infty)/R$ to the realised fraction. A quota over all rounds is met by solving with $\rho'=(\rho R-c_wT_w)/(R-T_w)$; both forms are reported, and the realised $\Pi_R$ is read from the logs.

\begin{figure}[t]
\centering
\begin{tikzpicture}[node distance=4mm and 6mm, every node/.style={font=\scriptsize}]
\node[term] (start) {round $r$};
\node[dec, below=of start] (due) {any tier due by (\ref{eq:due})?};
\node[proc, text width=36mm, below=of due] (bc) {broadcast the active tiers; sites train one epoch and return their slices};
\node[dec, below=of bc] (dp) {DP enabled?};
\node[proc, text width=36mm, below=of dp] (priv) {clip each delta to $C$ and sum;\\ add $\mathcal N(0,z^2C^2I)$;\\ divide by $qN$ (\ref{eq:dp})};
\node[proc, text width=36mm, below=of priv] (ag) {aggregate the active tiers (\ref{eq:agg}); $t_g\leftarrow r$};
\node[dec, below=of ag] (dp2) {DP enabled?};
\node[proc, text width=36mm, below=of dp2] (upd) {update $\Gh_g$ (\ref{eq:gammahat}); re-solve (\ref{eq:ip}) for $\mathbf K$};
\node[term, below=of upd] (next) {log bits per site; round $r+1$};
\node[proc, text width=17mm, right=7mm of bc] (local) {local-only round: zero payload, no release};
\draw[arr] (start) -- (due);
\draw[arr] (due) -- node[lab, right] {yes} (bc);
\draw[arr] (due.east) -| node[lab, above, pos=0.3] {no} (local);
\draw[arr] (local.south) |- (next.east);
\draw[arr] (bc) -- (dp);
\draw[arr] (dp) -- node[lab, right] {yes} (priv);
\draw[arr] (dp.west) -- ++(-9mm,0) node[lab, above, pos=0.5] {no} |- (ag.west);
\draw[arr] (priv) -- (ag);
\draw[arr] (ag) -- (dp2);
\draw[arr] (dp2) -- node[lab, right] {no} (upd);
\coordinate (k1) at ($(dp2.west)+(-9mm,0)$);
\coordinate (k2) at (k1 |- next.west);
\draw[line width=0.6pt] (dp2.west) -- node[lab, above] {yes} (k1);
\draw[line width=0.6pt] (k1) -- node[lab, sloped, midway] {planned $\mathbf K$} (k2);
\draw[arr] (k2) -- (next.west);
\draw[arr] (upd) -- (next);
\end{tikzpicture}
\caption{One server round. Under DP the schedule is fixed before training and the dispersion feedback is never computed from the private slices.}
\label{fig:round}
\end{figure}

\emph{Compression.} A compressor that emits a fraction $\kappa\in(0,1]$ of the float32 bits multiplies the upload, $B=R\,|\mathcal S|\,\Pi_R(\mathbf K)\,\kappa\,32\,d_{\mathrm{sh}}$ bits for $d_{\mathrm{sh}}$ shared parameters; it acts on each site's update relative to the broadcast, with error feedback, and the bits it emits are counted.

\subsection{Privacy: Threat Model, Mechanism and Eligible Schedules}
\label{sec:method-privacy}
\emph{Threat model.} Sites send their updates over TLS to an aggregator that is trusted to apply the mechanism and sees raw updates, as in central DP-FedAvg~\cite{mcmahan2018dp}. Consumers of the released models, the other sites among them, see every broadcast, and central DP protects the released aggregates against them. A passive network observer cannot read the payloads but sees when each site transmits and how many bytes, which is exactly the schedule; a schedule that depends on client data therefore leaks even when every payload is encrypted. Secure aggregation with distributed noise~\cite{kairouz2021ddg} would remove the trust in the aggregator but not this metadata, and poisoning is out of scope.

\emph{Mechanism.} Each round every site joins the cohort $\mathcal S_r$ independently with probability $q$; $q=1$, full participation, is the primary cross-silo setting and $q<1$ a sensitivity analysis. The server clips every participating site's returned delta and adds noise to the sum before dividing by the \emph{fixed} expected cohort size,
\begin{equation}
\tilde\Delta=\frac1{qN}\Big(\sum_{k\in\mathcal S_r}\Delta_k\min\big(1,\tfrac{C}{\|\Delta_k\|_2}\big)+\mathcal N(0,z^2C^2I)\Big),
\label{eq:dp}
\end{equation}
the fixed-denominator DP-FedAvg estimator~\cite{mcmahan2018dp}, so that adding or removing one site changes the sum by at most $C$ in $\ell_2$ and the denominator reveals nothing about the realised cohort; a planned release whose cohort is empty still adds noise. The mechanism is applied by the server to every method identically.

\emph{Accountant.} Over the $R_{\mathrm{tx}}$ planned releases the R\'enyi DP of the Poisson-subsampled Gaussian mechanism is composed at integer orders~\cite{mironov2019sgm},
\begin{align}
\varepsilon_\alpha&=\frac{R_{\mathrm{tx}}}{\alpha-1}\log\sum_{k=0}^{\alpha}\binom{\alpha}{k}(1-q)^{\alpha-k}q^ke^{\frac{k^2-k}{2z^2}},\notag\\
\varepsilon&=\min_\alpha\Big\{\varepsilon_\alpha+\frac{\log(1/\delta)}{\alpha-1}\Big\},
\label{eq:rdp}
\end{align}
with $\delta=10^{-5}$; for $q=1$ this reduces to $R_{\mathrm{tx}}\alpha/(2z^2)$~\cite{mironov2017}. The implementation reproduces the TensorFlow-Privacy MNIST reference ($\varepsilon\approx3.0$) and the $q=1$ closed form as unit tests. At $q=1$ the exact privacy of $R_{\mathrm{tx}}$ Gaussian releases is $\mu$-Gaussian DP with $\mu=\sqrt{R_{\mathrm{tx}}}/z$~\cite{dong2022gdp,balle2018}. Both accountants depend on $R_{\mathrm{tx}}$ only through $R_{\mathrm{tx}}/z^2$, so the exact one lowers every method's noise multiplier by the same factor, \num{gdpSigmaCutEightPct}--\num{gdpSigmaCutOnePct}\% for $\varepsilon^\star$ from 8 to 1. The comparison at matched $\varepsilon$ is therefore unchanged, and every reported $\varepsilon$ is an upper bound; $z$ is calibrated to each method's own release count by bisection on (\ref{eq:rdp}).

\emph{Which schedules are covered.} The accountant covers a release only if the decision to release does not depend on the private data. This holds for FedAvg, FedProx, FedPer, local training, a fixed period and the planner below, whose release rounds are fixed before training, and not for adaptive \tierfed{}, AdaComm, LAG or FedLAMA, whose schedules are functions of client updates or losses; the released code refuses to run them under DP. Data-dependent release decisions are a known pattern, which privacy filters and odometers~\cite{rogers2016odometer} and the sparse vector technique~\cite{lyu2017svt} charge for; the point here is that FL schedulers take such decisions without charging for them. DP-FTRL~\cite{kairouz2021ftrl} and matrix-factorisation mechanisms~\cite{choquette2023banded} suit full participation but also need data-independent release rounds; they are not evaluated.

\emph{Planning under a privacy budget (\tierfed{}-DP).} A planner that respects the accountant chooses $\mathbf K$ before training from public quantities only, $(\varepsilon^\star,\delta)$, $q$, $R$, $T_w$, $w_g$ and $\rho$. For each candidate it computes the release set of (\ref{eq:due}), the smallest $z(\mathbf K)$ that keeps those releases $(\varepsilon^\star,\delta)$-DP, and minimises
\begin{equation}
J(\mathbf K)=\frac{\sum_g w_gK_g^2}{K_{\max}^2}+\lambda\,\frac{\sum_{r}P_r(\mathbf K)\,z(\mathbf K)^2}{\sum_{r}P_r(\mathbf 1)\,z(\mathbf 1)^2},
\label{eq:jdp}
\end{equation}
the drift term of (\ref{eq:bound}) under a flat heterogeneity prior plus the injected noise energy, each normalised by its worst case, with $\lambda=1$. Fewer releases need less noise per release but drift more. Under the flat prior, $\Gamma_g\propto w_g$ and Proposition~\ref{prop:cuberoot} gives $\Psi=0$, so the planner can only choose equal intervals; it is evaluated as an ablation, not claimed as a contribution.
